\documentclass[11pt]{article}
\usepackage{fontspec}
\usepackage[margin=1.25in]{geometry}
\usepackage{amsmath,amssymb}
\usepackage[hidelinks]{hyperref}
\usepackage{setspace}
\setmainfont{TeX Gyre Pagella}
\setstretch{1.08}
\usepackage{titlesec}
\titleformat{\section}{\normalfont\Large\bfseries}{\thesection}{1em}{}
\titleformat{\subsection}{\normalfont\large\bfseries}{\thesubsection}{1em}{}

\title{\bfseries Access Before Storage: \\ A Continuation-Ontology Reading of Parametric Factuality}
\author{Flyxion \\ \small Independent Researcher}
\date{}

\begin{document}
\maketitle

\begin{abstract}
\noindent Ordinary factuality benchmarks ask whether a model's answer matches a ground truth, and when it does not, they attribute the failure to the model without asking which of several distinct things could have failed. A recent benchmark, WikiProfile, sharpens this by separating two conditions: whether a fact can be produced under a strongly primed, source-like prompt (``encoding''), and whether it can be produced under four differently phrased, minimally supported questions (``knowing''). The gap between these two conditions is real, measurable, and large even in frontier models. But the paper's language for the gap---shelves and keys, storage and retrieval---imports a container ontology that the behavioral evidence does not require and does not fully support. This essay develops an alternative: parameters as a compressed residue of training trajectories, and a produced answer as evidence of a locally admissible continuation from a particular prompt-state, not evidence of an object retrieved from a place. Read this way, the WikiProfile results become a single, useful worked example of a more general claim---that competence is not a property of a stored state considered alone, but a property of whether historically formed distinctions can be preserved, reactivated, reconstructed, and carried forward through an admissible continuation into a warranted output.
\end{abstract}

\section{The container ontology and its costs}

When a language model fails to produce a fact it was earlier shown to produce under a different prompt, the natural vocabulary is spatial. The fact was ``there'' and is now ``lost.'' It sits on a ``shelf'' that is either empty or blocked by a ``lost key.'' This vocabulary is not decorative. It carries a specific empirical commitment: that somewhere in the parameters there exists a discrete, addressable object corresponding to the proposition, and that the various prompts which do or do not produce the proposition are better or worse addresses pointing at that object.

The commitment is worth stating plainly because it is almost never tested directly. What is tested, in essentially every factuality benchmark including the most careful ones, is behavior under a finite set of prompts. A model is given a question, it produces or fails to produce a string, and the string is graded against a reference answer. Nothing in this procedure inspects the parameters. Nothing distinguishes a genuinely stored proposition from a locally reconstructible one, a partially overlapping association, or an artifact of how the prompt happens to constrain the space of plausible completions. The container ontology is an inference from behavior, not an observation of structure, and once it is stated as an inference rather than a finding, several things that looked settled become open again.

This essay is not a rebuttal of the empirical claim that models frequently fail to produce facts they can produce under other conditions. That claim is well supported and, if anything, undersold by how conservative the underlying experiments are. Nor is the container ontology an uncharitable reading forced onto a paper that only gestures at it in passing: the paper's own explanation of why thinking helps is framed as a search for ``systems in which the conditions for storing information diverge from the conditions for recalling it,'' which states storage-and-retrieval as the working hypothesis, not a decorative title-page metaphor. The essay proposes replacing that ontology with one built from admissible continuation, and shows what changes and what survives when the replacement is made. The WikiProfile benchmark \cite{calderon2026}, described in detail in the next section, serves throughout as the worked example against which the two ontologies can be compared, because it is unusually well instrumented: it distinguishes generation from recognition, direct phrasing from reverse phrasing, and unaided answering from answering after intermediate computation, in a single controlled design across thirteen models and roughly 4.5 million responses.

\section{The worked example: encoding, knowing, and the shape of the gap}

The benchmark in question, built around a collection of 2,150 Wikipedia-derived facts, tests each fact along two axes. The first axis, called encoding, asks whether the model can reproduce the fact under a prompt engineered to resemble the context in which the fact was probably encountered during training---a strongly primed, source-like condition. Success on either of two such prompts is sufficient to call the fact encoded. The second axis, called knowing, asks whether the model can reproduce the fact under four differently constructed questions: two seeking the object of the original proposition and two seeking its subject, each phrased two ways. Success on all four is required to call the fact known.

Two features of this design deserve attention before any interpretation is attempted. First, the two criteria are not symmetric tests of the same underlying property measured under easy and hard conditions; they are built from different logical quantifiers. Encoding is existential over a small, highly supportive set: it is satisfied if any one of two favorable prompts succeeds. Knowing is universal over a larger, less supportive set: it is satisfied only if every one of four minimally supported prompts succeeds. A fact needs one open door to count as encoded and cannot have any closed door to count as known. Some portion of the gap between the two rates is therefore already implied by the quantifier structure of the definitions, independent of any fact about what is or is not stored inside the model. This does not make the gap uninteresting---the paper's robustness checks show it survives across thresholds---but it changes what the size of the gap is allowed to mean. A large gap is guaranteed to be at least partly definitional; the empirical question is how much of it is more than that.

Second, the reported headline finding is that frontier models satisfy the encoding criterion for nearly all benchmark facts, on the order of 95 to 98 percent, while satisfying the knowing criterion for substantially fewer, with 25 to 33 percent of facts unrecoverable without intermediate computation and 11 to 12 percent unrecoverable even with it. Two conditions predict which facts fall into the gap: low popularity of the fact in the training distribution, and reverse phrasing, where the question asks for the subject of the original proposition rather than its object. Both effects are large and consistent across model families. Multiple-choice recognition of the correct answer is markedly easier than open generation of it, particularly in the reverse direction, where a model may be unable to generate the subject from the object but can reliably select it from four alternatives.

These are the results this essay treats as data. What is at stake is not whether the pattern exists but what description of the underlying system best explains it without smuggling in more structure than the behavior licenses.

\section{Parameters as residue, not container}

Consider what training actually does to a language model's parameters. It does not write propositions into addressed locations. It repeatedly adjusts a very large set of continuous weights so that, in aggregate and under whatever sampling of contexts the training data provided, certain continuations become more probable than others. The parameters at the end of this process are a compressed residue of an enormous number of local transitions, not a filing system. There is no reason, prior to specific evidence, to expect that residue to organize itself around the discrete, retrievable, direction-neutral objects that ordinary factual language treats as ``facts.''

This has a direct consequence for how to read the encoding test. A source-like prompt does not open a door onto a stored object; it partially reconstructs the conditions under which a particular trajectory was originally reinforced. If the reconstruction is close enough, the trained transition function completes it, and the correct string appears. Call this a reachable transition:
\[
x_{\text{prompt}} \longrightarrow x_{\text{answer}}.
\]
The object of interest is the transition itself---the admissible path from a specific starting state to a specific answer state---not a proposition sitting somewhere waiting to be pointed at. What the encoding test establishes is that at least one such path exists from at least one prompt. It does not establish that a path exists from every prompt, and it does not establish that the paths which do exist all pass through, or depend on, a single shared internal structure corresponding to the proposition in the ordinary sense.

This reframing explains the directionality result without any additional machinery. ``The band whose first gig was at the Boardwalk was Oasis'' and ``Oasis played its first gig at the Boardwalk'' are the same fact under a symmetric, propositional reading of knowledge. They are not the same trajectory. Autoregressive training builds directional continuations through text as it was actually written, and text is overwhelmingly written in the direction that introduces the more familiar entity first. The forward direction may have been traversed and reinforced many times; the reverse direction may have almost no direct support and must instead be assembled, if it can be assembled at all, from partial cues, co-occurrence, and elimination. Under a storage ontology this looks like the same fact being present and then mysteriously partly absent. Under a residue-and-trajectory ontology it is exactly what should be expected: two different paths through the same territory, built by two very different amounts of historical traffic.

The long-tail result is the same phenomenon at a different scale. A fact that was mentioned rarely in training leaves a thinner residue and fewer reinforced approach paths, so a minimally supportive prompt is less likely to land on one of them even though a strongly supportive, source-like prompt---closer to the one or two contexts where the fact actually appeared---still can. Popularity is not a proxy for whether the fact is stored; it is closer to a proxy for how many admissible entrances exist into the same underlying trajectory.

\section{What the encoding test actually shows}

It is worth being precise about what kind of claim the encoding test supports, because the paper's own language sometimes moves past it. The test shows that a fact survives at least one projection of the model into observable output under favorable conditions. Call this observational admissibility: there exists some prompt $q$ in the small, favorable set $E_f$ such that the model's output under $q$ matches the fact $f$,
\[
\exists\, q \in E_f : \pi(M, q) = f.
\]
This is a real and useful thing to establish. It is not the same as continuation admissibility---the claim that the fact is durably represented and would remain reachable from a wide, unprivileged range of future starting points. The paper's five-way profile of facts (encoding failure, recall failure despite thinking, recall recoverable only with thinking, direct recall, and inference without detected encoding) is built by combining an observational-admissibility test with a much stricter, closer-to-continuation test and treating the difference between them as informative about an internal storage state. The difference is informative, but about accessibility under specific conditions, not directly about storage.

The recognition experiments deserve the same care. Selecting the correct answer from four options demonstrates that the model can discriminate a supplied candidate from supplied alternatives---an admission decision made over a candidate set someone else constructed. It does not by itself demonstrate that the model could have generated that candidate on its own, and it does not demonstrate a symmetric, direction-independent representation of the relation between subject and object. The candidate has already crossed one boundary, from unconstrained generation into a bounded choice set, before the model does anything. The paper is careful about distractor quality and filtering, which rules out the crudest confounds, but it does not remove the structural fact that recognition and generation are different operations being asked to stand in for the same underlying property.

\section{Repair, not retrieval: what thinking does}

The paper's most striking secondary result is that inference-time computation---giving the model space to reason before answering---recovers a large fraction of facts that pass the encoding test but fail direct recall, on the order of 40 to 65 percent in thinking-optimized frontier models, against only 5 to 20 percent of facts that failed encoding altogether. The natural description offered is that thinking functions as a recall mechanism: it generates intermediate context that reactivates knowledge already present in the parameters. The same recovery shows up as a narrowing of the two gaps identified earlier: the popularity gap between long-tail and popular facts falls from a difference of 21.4 points to 12.5, and the directionality gap between reverse and direct questions falls from 9 points to 2. Both narrowings are consistent with repair in the sense argued for below, but neither is decisive on its own, for the reason given next.

This is a good description of one mechanism and an ambiguous description of the data as a whole, because the single category ``recall with thinking'' plausibly contains at least three distinct operations, and the experiments do not separate them. The first is genuine repair: intermediate computation reconstructs enough of the context that was present at encoding time to reconnect a poorly supported prompt to an already-reinforced trajectory, without introducing any new relation the model did not already have. The second is reconstruction: intermediate computation derives the answer from other, related propositions that are more directly accessible, arriving at the correct output through a different route than the one that would have produced it under source-like priming. The third, less discussed by the paper but visible in the reverse-question results, is closer to local search over a discriminable candidate set generated internally, which behaves more like recognition than like either repair or reconstruction.

These three operations have different implications. Repair is exactly the paper's intended story, and the evidence---that thinking gains concentrate on facts that already pass the encoding test, and that gains show up as increased consistency of correct answers rather than merely increased diversity of attempts---is genuinely consistent with it. But the same evidence is also consistent with reconstruction, since related propositions accessible via source-like priming would show the same concentration pattern without requiring that the target proposition itself be independently stored. Distinguishing the two is not a matter of designing a better prompt; it requires an intervention that can tell whether the same internal structure is doing the work in both the priming condition and the thinking condition, which is a different kind of experiment from anything in the paper.

\section{Sampling failure is not admissibility failure}

A further gap between what is measured and what is claimed concerns the treatment of unsuccessful attempts. Each question is sampled eight times, and a fact is scored as inaccessible under a given condition if none of the eight samples clears a grading threshold. This is a reasonable practical design, but it quietly converts a continuous production probability into a binary admissibility judgment. If the correct answer has a small but nonzero probability under the sampling distribution,
\[
P(f \mid q) \ll 1,
\]
eight samples will frequently miss it even though the transition is not, in the relevant sense, unavailable. The paper's language of facts that ``cannot be recalled'' describes a practical inaccessibility under a specific sampling regime and a specific number of attempts. It does not establish that the underlying transition function assigns the answer zero or near-zero weight, and it does not distinguish a genuinely excluded continuation from a weakly weighted one that a different decoding strategy, a longer sampling budget, or a slightly different prompt would surface. The authors themselves flag a version of this limitation, noting that with eight samples a correctness probability below roughly 0.1 is likely to yield zero correct draws, so that availability gains in the low-probability tail may be underestimated even though robustness at higher probabilities is not. That concession is narrower than the point made here: it treats the risk as an estimation problem confined to a low-probability tail, whereas the deeper issue is that any binary admissibility judgment drawn from a fixed number of samples collapses a continuous quantity, at every probability level, into present-or-absent. This matters because the distinction between exclusion and low weighting is exactly the distinction between two very different repair problems, addressed in the next section.

\section{Three kinds of repair deficit}

The paper's five-profile taxonomy is a genuine improvement over single-score accuracy, but it groups together failures that plausibly have different causes and would respond to different interventions. It is useful to separate at least three.

A structural deficit exists when no admissible transition connects a class of prompts to the correct answer at all---nothing in the trained function supports the path, regardless of how much computation or context is supplied. This is the closest behavioral approximation to genuine non-encoding, though even here the caveat from the previous section applies: absence under eight samples is evidence for, not proof of, a structural deficit.

An energetic deficit exists when an admissible transition does exist but cannot be executed within the compute, context length, or decoding budget available at inference time. Reverse questions plausibly fall here more often than the paper's framing suggests: the transition is not missing, it is expensive to traverse without an intermediate step that reconstructs the missing directional support, which is exactly what thinking sometimes supplies.

A kinetic deficit exists when recovery is possible and would eventually succeed but is too slow, in the sense of requiring more intermediate steps or more sampling than a fixed budget allows, relative to whatever latency or cost constraint the task imposes. This is distinct from an energetic deficit in that it is not a hard budget failure but a soft tradeoff, and it is the category most directly relevant to the paper's own observation that models struggle to decide, without being told, when a question warrants the extra cost of thinking at all.

None of these three deficits requires positing an object that is present or absent on a shelf. All three are properties of the transition structure and the resources available to traverse it, which is precisely the vocabulary a residue-and-trajectory ontology provides and a storage ontology does not.

\section{Non-collapse: four things that are not the same}

The overall discipline this essay is arguing for can be stated as a short non-collapse chain. Successful completion of a prompt does not entail that a discrete proposition was stored. A stored residue, even granting some structure corresponding to it, does not entail robust knowledge in the sense of being reachable from a wide range of unprivileged starting points. Robust reachability, in turn, does not entail that an emitted answer constitutes a warranted commitment---a claim the system stands behind independent of which particular prompt happened to elicit it:
\begin{center}
\small
successful completion \;$\not\Rightarrow$\; stored proposition \;$\not\Rightarrow$\; robust knowledge

\vspace{2pt}
robust knowledge \;$\not\Rightarrow$\; warranted commitment.
\end{center}
The WikiProfile design is valuable precisely because it separates the middle two terms of this chain more carefully than almost any comparable benchmark: its encoding/knowing split is functionally a measurement of the gap between the first and third terms. Where it remains vulnerable to the container metaphor is in occasionally treating success at the first term as though it had already established the second, and in treating success at the third term, when it happens, as though it settled the fourth. A fact recalled under a strongly primed, one-off condition is not yet a commitment the system would defend under a different framing, a different questioner, or a different context---the same concern developed at length, for historical records and provenance claims, in a companion essay on declared versus witnessed history, and which applies with equal force to a model's momentary factual output.

\section{What would count as stronger evidence}

Everything above is compatible with the storage vocabulary being roughly right; the point is that the WikiProfile experiments, on their own, cannot distinguish it from the alternative. Settling the question would require moving from behavioral production to causal structure. A minimal program would identify a candidate internal representation associated with a specific fact, using whatever probing or activation-analysis method is available; perturb or suppress that representation directly rather than varying the prompt; test whether the perturbation produces selective loss of the fact across the full range of conditions---source-like priming, direct questions, reverse questions, and recognition---rather than loss under only one of them; restore the representation and confirm recovery under the same range of conditions; and, critically, check whether an intact model answering correctly via thinking is using the same representation that source-like priming activates, or is instead constructing the answer through a different, derivational route that happens to converge on the same output string. Only the last step can separate repair from reconstruction in the sense argued for above. None of this is a criticism of the paper for not doing it---behavioral benchmarking at the scale attempted here is already a substantial undertaking, and the interventionist program is a different kind of project with different costs. It is a specification of the gap between what has been shown and what the shelves-and-keys language suggests has been shown.

\section{A categorical and sheaf-theoretic restatement}

The argument so far has been made in prose, with a handful of set-theoretic definitions borrowed directly from the paper's own notation. It is worth pausing to give the same argument a slightly more structural restatement, not because the phenomenon requires heavier machinery to be true, but because the machinery makes precise something the prose can only gesture at: exactly what kind of gap the encoding/knowing split is measuring, and why no amount of local success closes it.

Fix a fact $f$ and let $Q_f$ denote the full set of contexts under which the paper tests it: the two source-like priming prompts, the four generative questions, and the recognition variants. Treat $Q_f$ as an index category with no nontrivial morphisms between contexts---the point is not that one context factors through another, only that $Q_f$ can be covered by distinguished subsets, exactly as an open set is covered by smaller open sets in a topological space. The relevant subsets are the encoding cover $E_f \subset Q_f$ and the four-element knowing cover $K_f \subset Q_f$, with $E_f$ properly contained in $Q_f$ and $K_f$ properly contained in $Q_f$, and the two covers only partially overlapping.

Define a presheaf of successful completions $\mathcal{F}_f : Q_f^{\mathrm{op}} \to \mathbf{Set}$, where $\mathcal{F}_f(q)$ is the one-point set $\{\ast\}$ if the model produces $f$ under context $q$ and the empty set otherwise. This is a boolean-valued presheaf, and asking whether it has a section over a cover is exactly asking whether every context in that cover is answered correctly. The encoding test asks whether $\mathcal{F}_f$ has a section over some single object of $E_f$---an existential, local claim, satisfied as soon as one stalk is nonempty. The knowing test asks whether $\mathcal{F}_f$ has a global section over the whole cover $K_f$---every stalk in $K_f$ nonempty, and, because the presheaf is boolean-valued, automatically compatible wherever sections overlap, so the only obstruction to gluing is a stalk that is simply empty.

This is the precise sense in which the paper's headline gap is a locality phenomenon rather than a storage phenomenon. A presheaf can have nonempty stalks at isolated points of its domain while having no global section over a larger cover, without there being anything wrong with the presheaf considered as an object; this is the ordinary situation of a function defined only on part of a space. Calling the domain where $\mathcal{F}_f$ is nonempty a ``stored fact'' is analogous to calling a function's domain of definition a location inside its codomain. The reversal curse \cite{berglund2024}, on this reading, is the observation that $\mathcal{F}_f$ is frequently nonempty on the direct-question objects of $Q_f$ and empty on the reverse-question objects, which is an ordinary failure of a section to extend, not evidence that two different propositions are involved.

Thinking's effect can then be described as changing the presheaf itself rather than merely improving access to a fixed one. Let $\mathcal{F}_f^{\mathrm{think}}$ be the analogous presheaf evaluated with intermediate computation permitted before each answer. The paper's finding that thinking recovers a large share of encoded-but-not-known facts is the claim that $\mathcal{F}_f^{\mathrm{think}}$ has strictly more nonempty stalks than $\mathcal{F}_f$ over $K_f$, concentrated on facts where $\mathcal{F}_f$ was already nonempty somewhere in $E_f$. Repair, in the sense argued for in an earlier section, is the special case where the added stalks arise from a map that factors through the same stalk that made $\mathcal{F}_f$ nonempty on $E_f$ in the first place---an actual morphism connecting the primed context to the previously empty context. Reconstruction is the case where the added stalks arise without any such factoring, from a different route entirely that happens to land on the same value. Nothing in the paper's design distinguishes these two cases, because nothing in the paper's design tracks factoring; it only tracks whether stalks are empty or not. This is the categorical restatement of the causal-evidence gap identified above: the presheaf formalism makes visible exactly the information a purely behavioral test cannot access, namely whether an extension of a section factors through the original section or is independently constructed.

A brief note on types is worth adding here, since the same distinction recurs in a more familiar form under the Curry--Howard correspondence between proofs and programs. Producing an answer under a specific prompt is naturally read as a typing judgment relative to a context: $\Gamma_q \vdash t_f : \mathrm{Fact}(f)$, where $\Gamma_q$ packages whatever the prompt $q$ supplies and $t_f$ is the term the model outputs. The encoding test establishes that some such judgment is derivable for at least one favorable $\Gamma_q$. What the knowing test is really asking is whether a term of type $\mathrm{Fact}(f)$ is derivable in a context-independent sense---closer to asking for a closed term, one that does not depend on which particular $\Gamma_q$ happens to be supplied. A term derivable only under one specific, richly informative context is not thereby derivable under the empty context or under an unfavorable one, and the gap between an open term with free dependence on $\Gamma_q$ and a genuinely closed term is precisely the gap between encoding and knowing restated in typing vocabulary. This is also why recognition is a different judgment from generation rather than a weaker version of the same one: recognition supplies the candidate term itself as part of the context and asks only for a judgment of type equality against the alternatives, which is a strictly easier problem than deriving the term from $\Gamma_q$ alone, and no amount of success at the easier problem entails success at the harder one.

\section{Conclusion}

The value of the WikiProfile results does not depend on accepting that models contain facts which are sometimes hard to reach. It depends only on the much better supported claim that factual production is radically conditional: the same model, given the same underlying association, will produce it, fail to produce it, or produce it only after intermediate computation, depending on the direction, popularity, and phrasing of the question. That conditionality is the real discovery, and it survives intact under the alternative ontology proposed here. What does not survive is the inference from ``conditional production'' to ``discrete stored object, sometimes hard to find.'' The more defensible description is that a model's parameters are a compressed residue of many historical training trajectories; that a produced answer is evidence of one admissible continuation from one particular prompt-state, not evidence of retrieval from a place; that thinking, when it helps, is best understood as a repair operator that sometimes reconnects a weakly supported prompt to an already-reinforced trajectory and sometimes reconstructs the answer by an entirely different route; and that none of the available behavioral evidence, careful as it is, yet distinguishes those two cases. Factual competence, on this reading, is not a property a state possesses. It is a property of whether a historically formed distinction can be preserved, reactivated, reconstructed, and carried through an admissible continuation into an output the system would actually stand behind.

\begin{thebibliography}{9}

\bibitem{calderon2026}
N. Calderon, E. Ben-David, Z. Gekhman, E. Ofek, and G. Yona.
Empty shelves or lost keys? Recall is the bottleneck for parametric factuality.
In \textit{Proceedings of the 43rd International Conference on Machine Learning}, ICML 2026.
arXiv:2602.14080.

\bibitem{berglund2024}
L. Berglund, M. Tong, M. Kaufmann, M. Balesni, A. C. Stickland, T. Korbak, and O. Evans.
The reversal curse: LLMs trained on ``A is B'' fail to learn ``B is A''.
In \textit{The Twelfth International Conference on Learning Representations}, ICLR 2024.

\end{thebibliography}

\end{document}

